\chapter{Background}\label{ch:background}

\section{Design Patterns}\label{sec:design-patterns}

Design patterns can be defined as a common design structure useful for creating reusable object-oriented designs. They
are characterized by four essential elements~\cite{gamma1995designpatterns}:
\begin{itemize}
    \item The name provides a concise way to refer to the pattern and captures the intended application domain.
    \item The problem describes the considered design issues and the context.
    \item The solution is expressed as the elements that make up the design, their relationships, responsibilities, and collaborations.
    \item The consequences of applying the pattern and its and trade-offs.
\end{itemize}

This work starts from the original dataset of the LLM Design Pattern Refactoring project~\cite{zisa2024llmdesignpatternrefactoring},
which contained three examples for each of the following design patterns: Strategy, State, Decorator, Adapter, and
Factory Method. Each example was provided in two versions: a refactored version implementing the corresponding design
pattern and a baseline version, where the pattern had been removed.

Starting from this dataset, this project extends the set of considered design patterns by introducing five additional
patterns, briefly described in the following paragraphs.

The first added pattern is Abstract Factory, which provides an interface for creating families of related or dependent
objects without specifying their concrete classes~\cite{gamma1995designpatterns}.
\begin{figure}[htbp]
    \centering
    \includegraphics[width=0.8\textwidth]{figures/design_patterns_diagrams/abstract_factory}
    \caption{Abstract Factory pattern structure}
    \label{fig:abstract-factory-pattern-structure}
\end{figure}
\FloatBarrier

Builder pattern separates the construction of a complex object from its representation so that the same construction
process can create different representations~\cite{gamma1995designpatterns}.
\begin{figure}[htbp]
    \centering
    \includegraphics[width=0.4\textwidth]{figures/design_patterns_diagrams/builder}
    \caption{Builder pattern structure}
    \label{fig:builder-pattern-structure}
\end{figure}
\FloatBarrier

Proxy pattern provides a surrogate or placeholder for another object to control access to it~\cite{gamma1995designpatterns}.
\begin{figure}[htbp]
    \centering
    \includegraphics[width=0.4\textwidth]{figures/design_patterns_diagrams/proxy}
    \caption{Proxy pattern structure}
    \label{fig:proxy-pattern-structure}
\end{figure}
\FloatBarrier

Command pattern encapsulates a request as an object, thereby letting you parameterize clients with different requests,
queue or log requests, and support undoable operations~\cite{gamma1995designpatterns}.
\begin{figure}[htbp]
    \centering
    \includegraphics[width=0.7\textwidth]{figures/design_patterns_diagrams/command}
    \caption{Command pattern structure}
    \label{fig:command-pattern-structure}
\end{figure}
\FloatBarrier

Template Method pattern defines the skeleton of an algorithm in an operation, deferring some steps to subclasses so that
they can redefine certain steps without changing the algorithm’s structure~\cite{gamma1995designpatterns}.
\begin{figure}[htbp]
    \centering
    \includegraphics[width=0.4\textwidth]{figures/design_patterns_diagrams/template_method}
    \caption{Template Method pattern structure}
    \label{fig:template-method-pattern-structure}
\end{figure}
\FloatBarrier

Each design pattern focuses on a particular object-oriented design problem or issue. It describes when it applies,
whether it can be applied in view of other design constraints, and the consequences or trade-offs of its use.
Design patterns are classified into three main categories, depending on their purpose:
\begin{itemize}
    \item Creational patterns concern the process of object creation
    \item Structural patterns deal with the composition of classes or objects.
    \item Behavioral patterns characterize the ways in which classes or objects interact and distribute responsibility.
\end{itemize}

In this work, the considered design patterns cover all three categories. Strategy, State, Command, and Template Method
belong to behavioral patterns; Decorator, Adapter, and Proxy belong to structural patterns; while Factory Method,
Abstract Factory, and Builder belong to creational patterns.

\section{Large Language Models for Code}\label{sec:large-language-models-for-code}

\section{Structured outputs}\label{sec:structured-outputs}

\section{Ranking-based evaluation}\label{sec:ranking-based-evaluation}
